% Generated draft: requires mathematical and source review.
\documentclass[11pt]{article}
\usepackage[margin=25mm]{geometry}
\usepackage{fontspec}
\setmainfont{Latin Modern Roman}
\usepackage{amsmath,amssymb,mathtools,unicode-math}
\setmathfont{Latin Modern Math}
\usepackage{enumitem,needspace,xurl,hyperref}
\setlength{\parindent}{0pt}
\setlength{\parskip}{5pt}
\setlength{\emergencystretch}{3em}
\newcommand{\N}{\mathbb{N}}
\newcommand{\Z}{\mathbb{Z}}
\newcommand{\Q}{\mathbb{Q}}
\newcommand{\R}{\mathbb{R}}
\newcommand{\C}{\mathbb{C}}
\newcommand{\F}{\mathbb{F}}
\newcommand{\A}{\mathbb{A}}
\newcommand{\PP}{\mathbb{P}}
\newcommand{\E}{\mathbb{E}}
\newcommand{\eps}{\varepsilon}
\newcommand{\dd}{\,\mathrm{d}}
\begin{document}
\section*{UnsolvedMath ID: 10800013}
\textbf{Source problem code:} catalogue entry
\subsection{Definitions and mathematical statement}
\emph{Problem (Victor Vassiliev, Problem 5C).} Give an upper bound for the function \(T\mapsto F\) defined in the omitted source discussion. Writing its dependence schematically as \(F(T)\), the requested bound has the form
\[
 F(T)\leq B(T)\qquad\text{for all admissible }T,
\]
with any other parameters and their quantifiers specified in the intended setting. The supplied statement does not define \(T\), \(F\), their domain or range, or the desired kind of bound. Thus this inequality records the request only; it does not introduce a defined function or a particular conjectural rate.

The supplied research summary places the function in the setting of the discriminant stratification of a space denoted \(P(d;N)\). It mentions classifications for \(N=2\) and small \(d\) as possible sources of data. The record does not define \(P(d;N)\), the relation of its strata or components to \(T,F\), whether \(d,N\) are fixed, or what uniformity is required. Those definitions are essential before an upper-bound problem can be made quantitative. No polynomial, exponential, effective, or asymptotic bound is silently selected.
\subsection{Short English statement}
Find an upper bound for the source's function \(T\mapsto F\). Its variables and their relationship to the polynomial-space discriminant are omitted, so the draft records the bound request without inventing a function or growth rate.
\subsection{Sources}
\begin{itemize}
\item Local UnsolvedMath registry, permanent ID 10800013, source code catalogue entry; proposer Victor Vassiliev, Problem 5C.
\item Vassiliev, \emph{A Few Problems on Monodromy and Discriminants}, Problem 5C. Supplied URL: \url{https://amj.math.stonybrook.edu/html-articles/Files-2015-2024/15-11/}; recorded access 2026-07-29; extraction html-structure. The article contains eighteen explicitly labeled theorem-style problem environments.
\item The supplied research summary mentions Vassiliev's degree-4 Morse polynomial isotopy classification, \texttt{arXiv:2311.11113}, and ``the 64 sextics,'' \texttt{arXiv:1703.01660}, as providing implicit component data for \(N=2\) and small \(d\). No additional bibliographic details or verified identification of the relevant data are supplied.
\item Supplied status open and OPEN-TRIAGE. The summary says ``Open (as of Aug 2026),'' reports no published upper bound or relevant arXiv hits found, and says classifications might allow computation in some cases. These are supplied dated search assessments, not an independently verified absence of results. Current primary-literature status and redistribution terms for the publicly accessible author/publisher copy remain NEEDS\_REVIEW.
\end{itemize}
\end{document}
